\documentclass[11pt]{article}
\usepackage[a4paper,margin=27mm]{geometry}
\usepackage{amsmath,amssymb,amsthm,hyperref}
\newtheorem{proposition}{Proposition}
\title{A short quotient argument for one matrix quartet}
\author{Human Compression}
\date{}
\begin{document}
\maketitle
Write $\langle a,b,c,d\rangle$ for the subgroup generated by the supplied
matrices over $\mathbb F_{11}$. Matrix products act on rows from the right.
The computational certificate verifies the equations in the following
proposition. The finite presentation theorem is borrowed from
Hartley and Leemans~\cite[Theorem~3.2]{HL}; it is not reproved here.

\begin{proposition}
Suppose that $a,b,c,d$ are elements of a group, that each is an involution,
and that
\[
 ac=ca,\qquad ad=da,\qquad bd=db,\qquad
 (ab)^5=(bc)^3=(cd)^5=(abc)^5=1.
\]
Then $\lvert\langle a,b,c,d\rangle\rvert\leq600415200$.
In particular this group is not isomorphic to $G_2(11)$.
\end{proposition}
\begin{proof}
Let $W$ be the group with generators $s_0,s_1,s_2,s_3$, the Coxeter
relations of type $[5,3,5]$, and the additional relation
$(s_1s_2s_3)^5=1$. The cited theorem identifies this literal presented
group with $J_1\times\operatorname{PSL}_2(19)$, of order
$175560\cdot3420=600415200$.

Send $(s_0,s_1,s_2,s_3)$ to $(d,c,b,a)$. The involution and distant
commutation relations hold by hypothesis. The three adjacent products
are the inverses of $cd,bc,ab$, so their required powers are the identity.
The added word maps to $cba=(abc)^{-1}$, whose fifth power is also the
identity. The universal property of a group presentation therefore gives
a surjection $W\longrightarrow\langle a,b,c,d\rangle$.
The order of the latter group divides $|W|$, giving the claimed bound.
Finally,
\[
 |G_2(11)|=11^6(11^6-1)(11^2-1)=376611192619200>600415200.
\]
\end{proof}

The supplied tuple is identified with original F004, job20, position200
by the existing source and frame records. Its consecutive intersection
property is a separate earlier certificate. Neither that fact nor an
abstract recognition of either end subgroup is required in the argument
above. The application supplies a nongeneration proof for this tuple;
it is not a newly proved open conjecture.

\begin{thebibliography}{1}
\bibitem{HL} M. I. Hartley and D. Leemans,
\emph{Quotients of a Universal Locally Projective Polytope of type
$\{5,3,5\}$}, arXiv:math/0307075v1, 2003.
\url{https://arxiv.org/pdf/math/0307075v1}.
\end{thebibliography}
\end{document}
